\section{Introducción}

La producción audiovisual en directo ha experimentado en los últimos años una profunda
transformación, impulsada por la democratización del hardware de bajo coste y el
crecimiento exponencial del \textit{streaming} en Internet. Eventos técnicos de referencia como
el Chaos Communication Congress (CCC), organizado por el Chaos Computer Club en Alemania,
han puesto de manifiesto la necesidad de disponer de herramientas de realización
profesional que sean al mismo tiempo económicas, fiables y auditables.

En este contexto surge Voctomix, desarrollado por el equipo C3VOC (\textit{Chaos Communication
Congress Video Operations Crew}). Frente a soluciones previas como \textit{dvswitch}, limitada
en resolución y escalabilidad, o herramientas comerciales como vMix o OBS Studio,
diseñadas para un único operador en un único equipo, Voctomix adopta una arquitectura
cliente-servidor desacoplada que separa el motor de mezcla multimedia de la interfaz
de control, permitiendo distribuir la carga de procesamiento en distintos nodos de red.

El presente Trabajo de Fin de Grado toma como punto de partida la versión original
de Voctomix y propone una serie de extensiones orientadas a su uso en producciones
en directo de mediana escala: rotulación dinámica sobre el vídeo, gestión avanzada
del blanking de stream, composites avanzados de vídeo, telemetría en tiempo real
mediante un broker de mensajería, y la contenerización completa del sistema mediante
Docker y Docker Compose. El resultado es Voctomix 2.0, una plataforma más versátil,
desplegable y adaptada a flujos de trabajo profesionales.

Un elemento diferenciador de este proyecto es su orientación hacia un \textbf{caso de uso
real}: el sistema desarrollado está previsto para su utilización por el grupo de
investigación \textbf{CyberNemo} de la Universidad Politécnica de Madrid (UPM). Este
grupo necesita una solución de producción audiovisual remota que permita retransmitir
sus actividades, seminarios y presentaciones con calidad profesional y sin depender de
infraestructura hardware costosa. Voctomix 2.0 responde directamente a estas necesidades,
proporcionando un entorno reproducible y fácilmente desplegable que puede ponerse en
funcionamiento con un único comando (\texttt{docker compose up}) en cualquier equipo
Linux estándar.

% --- Figura sugerida ---
% \begin{figure}[H]
%   \centering
%   \includegraphics[width=0.85\textwidth]{figuras/gui_voctomix2.png}
%   \caption{Interfaz gráfica de Voctomix 2.0 en funcionamiento.}
%   \label{fig:gui_voctomix2}
% \end{figure}

% --- Figura sugerida ---
% \begin{figure}[H]
%   \centering
%   \includegraphics[width=0.85\textwidth]{figuras/diagrama_contexto.png}
%   \caption{Diagrama de contexto: Voctomix 2.0 en el entorno de producción del grupo CyberNemo (UPM).}
%   \label{fig:diagrama_contexto}
% \end{figure}
